\section{Proofs}
\label{app:proofs}

Throughout the appendix, \cref{ass:model} holds for every site that
appears, including candidate sites unless the text says otherwise, and
the eigenvalue condition of \cref{sec:setting} holds for the pool and,
where an augmented pool appears, for the augmented pool. We write
$\Pi_V$ and $\Pi^\perp$ for the orthogonal projections onto $\Vspan$ and
$\Vspan^\perp$, and $\beta^\circ=(\beta_0^\circ,\beta_1^\circ)$ for the
true coefficient vector of \eqref{eq:response-model}. Recall from
\eqref{eq:span} that $\Vspan$ is the span of the centered site profiles
and that $\affhull(\pool)=c_1+\Vspan$.

Two facts from \citet{boyarsky2026} are used repeatedly. First, the
normal equations: under \cref{ass:model}, for any pooled source law
whose sites obey \eqref{eq:response-model} with the same coefficients,
\begin{equation}
  \cov(C,\pseudo{a}\mid X)=\cov(C\mid X)\,\beta_a^\circ
  \quad\text{almost surely, so}\quad
  \dmom{a}=\CovC\beta_a^\circ.
  \label{eq:app-normal}
\end{equation}
Second, under the eigenvalue condition, $\range(\CovC)=\Vspan$ and
$\kernel(\CovC)=\Vspan^\perp$, so the solution set of the equality
system in \eqref{eq:feasible-set} is
$\{\beta_a:\ \CovC\beta_a=\dmom{a}\}
=\beta^\circ_a+\Vspan^\perp$.

\subsection{Proof of \texorpdfstring{\cref{lem:loading}}{Lemma 1}}
\label{app:proof-loading}

\begin{proof}
Write $p_s(x)=\pr_\sP(C=c_s\mid X=x)$, so
$\muC(x)=\sum_{s=1}^S p_s(x)\,c_s$ with $p_s(x)\geq0$ and
$\sum_s p_s(x)=1$ for $\sP$-almost every $x$, hence for $\tP$-almost
every $x$ by overlap. Then
$q:=\E_{\tP}[\muC(X)]=\sum_s \bar w_s c_s$ with $\bar w_s=\E_{\tP}[p_s(X)]\geq0$
and $\sum_s\bar w_s=1$, so $q\in\convhull\{c_1,\ldots,c_S\}\subseteq
\affhull(\pool)$. Since $\affhull(\pool)=c_1+\Vspan$, we have
$q-c_1\in\Vspan$ and therefore
\[
  \Pi^\perp\ellvec=\Pi^\perp(\ct-q)
  =\Pi^\perp(\ct-c_1)-\Pi^\perp(q-c_1)=\Pi^\perp(\ct-c_1).
\]
The distance from a point to the affine set $c_1+\Vspan$ is the norm of
the component of the offset orthogonal to $\Vspan$, which gives
\eqref{eq:loading-distance}. Finally, under the eigenvalue condition
$\range(\CovC)=\Vspan$, so $\ellvec\in\range(\CovC)$ if and only if
$\Pi^\perp\ellvec=0$, if and only if $\Pi^\perp(\ct-c_1)=0$, if and only
if $\ct\in c_1+\Vspan=\affhull(\pool)$.
\end{proof}

\subsection{Proof of \texorpdfstring{\cref{prop:dichotomy}}{Proposition 1}}
\label{app:proof-dichotomy}

\begin{proof}
With $\restrset=\R^{2p}$ the feasible set is the product of the two
equality solution sets, $\feasset=\prod_{a\in\{0,1\}}
(\beta^\circ_a+\Vspan^\perp)$, which contains $\beta^\circ$. For
$\beta_a=\beta^\circ_a+z_a$ with $z_a\in\Vspan^\perp$,
\[
  \ellvec^\top(\beta_1-\beta_0)
  =\ellvec^\top(\beta^\circ_1-\beta^\circ_0)
  +(\Pi^\perp\ellvec)^\top(z_1-z_0),
\]
because $z_1-z_0\in\Vspan^\perp$. If $\ct\in\affhull(\pool)$, then
$\Pi^\perp\ellvec=0$ by \cref{lem:loading}, the objective is constant on
$\feasset$, and both endpoints equal $\psi_T$, so the width is zero. If
$\ct\notin\affhull(\pool)$, then $\Pi^\perp\ellvec\neq0$. Taking
$z_1=t\,\Pi^\perp\ellvec$ and $z_0=0$ with $t\to\pm\infty$ drives the
objective to $\pm\infty$, so the width is infinite.
\end{proof}

\subsection{Proof of \texorpdfstring{\cref{prop:hull}}{Proposition 2}}
\label{app:proof-hull}

\begin{proof}
Write $u=\Pi^\perp(\cand{k}-c_1)$, so
$\cand{k}-c_1=\Pi_V(\cand{k}-c_1)+u$ and
\[
  \Vspan'=\Vspan+\operatorname{span}\{\cand{k}-c_1\}
  =\Vspan+\operatorname{span}\{u\}.
\]
If $u=0$ the two spans coincide. If $u\neq0$ the sum is orthogonal and
$\dim\Vspan'=r+1$; moreover $u=0$ exactly when
$\cand{k}\in c_1+\Vspan=\affhull(\pool)$, which proves the equivalences.
For the complement, $x\in\Vspan'^{\perp}$ if and only if $x\perp\Vspan$
and $x\perp u$, that is $\Vspan'^\perp=\Vspan^\perp\cap\{u\}^\perp$.
Finally, every residual $C-\muC'(X)$ of the augmented pool lies in
$\Vspan'$, because $\muC'(x)$ is a convex combination of the augmented
profiles, so $\range(\CovC')\subseteq\Vspan'$ always, and the augmented
eigenvalue condition states that
$\Bspan'^\top\CovC'\Bspan'$ is positive definite, which forces
$\range(\CovC')=\Vspan'$ and, by symmetry of $\CovC'$,
$\kernel(\CovC')=\Vspan'^{\perp}$.
\end{proof}

\subsection{Proof of \texorpdfstring{\cref{prop:distance}}{Proposition 3}}
\label{app:proof-distance}

\begin{proof}
Only the shift $\shift=\beta_1-\beta_0$ enters the objective and the
restriction, and the pair of equality systems constrains it through
$\CovC\shift=\dmom{1}-\dmom{0}$: given any such $\shift$, choosing
$\beta_0$ with $\CovC\beta_0=\dmom{0}$ and $\beta_1=\beta_0+\shift$
satisfies both systems. The solution set for the shift is
$\bar\shift+\Vspan^\perp$, where $\bar\shift$ is the minimum-norm
solution and lies in $\Vspan$. For $\shift=\bar\shift+\nu$ with
$\nu\in\Vspan^\perp$, orthogonality gives
$\norm{\shift}_2^2=\norm{\bar\shift}_2^2+\norm{\nu}_2^2$, so the cap
$\norm{\shift}_2\leq\tau$ is exactly
$\norm{\nu}_2\leq\rho_\tau:=(\tau^2-\norm{\bar\shift}_2^2)^{1/2}$, and
the objective is
$\ellvec^\top\shift=\ellvec^\top\bar\shift+(\Pi^\perp\ellvec)^\top\nu$.
By the Cauchy--Schwarz inequality, the supremum and infimum of
$(\Pi^\perp\ellvec)^\top\nu$ over the ball of radius $\rho_\tau$ in
$\Vspan^\perp$ are $\pm\rho_\tau\norm{\Pi^\perp\ellvec}_2$, attained at
$\nu=\pm\rho_\tau\,\Pi^\perp\ellvec/\norm{\Pi^\perp\ellvec}_2$ when
$\Pi^\perp\ellvec\neq0$ and trivially when it is zero. The width is
therefore $2\rho_\tau\norm{\Pi^\perp\ellvec}_2$, and
\cref{lem:loading} identifies the last factor with
$\dist(\ct,\affhull(\pool))$.
\end{proof}

\subsection{Proof of \texorpdfstring{\cref{prop:refinement}}{Proposition 4}}
\label{app:proof-refinement}

\begin{proof}
Apply \eqref{eq:app-normal} to the augmented pool: because the candidate
site obeys \eqref{eq:response-model} with the same coefficients,
$\dmom{a}'=\CovC'\beta^\circ_a$, so the augmented equality solution set
is $\beta^\circ_a+\kernel(\CovC')$, and by \cref{prop:hull} this kernel
is $\Vspan^\perp\cap\{\newdir{k}\}^\perp$ in the hull-expanding case and
$\Vspan^\perp$ in the rank-preserving case.

In the rank-preserving case the solution sets coincide with the pool's,
so $\feasset'=\feasset$ after intersecting with $\restrset$. In the
hull-expanding case, for $\beta_a=\beta^\circ_a+z_a$ with
$z_a\in\Vspan^\perp$ we have
$\newdir{k}^\top\beta_a=\newdir{k}^\top\beta^\circ_a
+\newdir{k}^\top z_a$, and $z_a\in\{\newdir{k}\}^\perp$ exactly when
$\newdir{k}^\top\beta_a=\newdir{k}^\top\beta^\circ_a=\reveal{a}{k}$.
Hence
\[
  \beta^\circ_a+\kernel(\CovC')
  =\bigl(\beta^\circ_a+\Vspan^\perp\bigr)
  \cap\bigl\{\beta_a:\ \newdir{k}^\top\beta_a=\reveal{a}{k}\bigr\},
\]
and intersecting both arms' sets with $\restrset$ gives
\eqref{eq:refinement}. The scalar $\reveal{a}{k}$ is a linear functional
of $\Bspan'^\top\beta^\circ_a$, which the augmented pool identifies
through its normal equations because $\newdir{k}\in\Vspan'$; it is not
identified by the pool alone because moving $\beta_a$ along
$\newdir{k}\in\Vspan^\perp$ leaves every pool moment unchanged.

For the last claim, compare the augmented-pool endpoint programs, which
use $(\bT',\ellvec')$, with the base-objective programs over
$\feasset'$. Write $q=\E_{\tP}[\muC(X)]$ and $q'=\E_{\tP}[\muC'(X)]$ for
the two mixture points. As in the proof of \cref{lem:loading}, $q$ lies
in $\affhull(\pool)$ and $q'$ in $\affhull(\pool\cup\{k\})$, so
$\ellvec-\ellvec'=q'-q\in\Vspan'$, and for every $\beta\in\feasset'$,
writing $\shift=\beta_1-\beta_0=\shift^\circ+\nu$ with
$\nu\in\Vspan'^\perp$,
\[
  (\ellvec-\ellvec')^\top\shift
  =(\ellvec-\ellvec')^\top\shift^\circ
\]
is constant on $\feasset'$. Both decompositions represent the same
target effect, $\psi_T=\bT+\ellvec^\top\shift^\circ
=\bT'+\ellvec'^\top\shift^\circ$, so
$\bT'-\bT=(\ellvec-\ellvec')^\top\shift^\circ$, and for every
$\beta\in\feasset'$,
\[
  \bT'+\ellvec'^\top\shift
  =\bT+(\ellvec-\ellvec')^\top\shift^\circ
  +\ellvec^\top\shift-(\ellvec-\ellvec')^\top\shift
  =\bT+\ellvec^\top\shift .
\]
The two programs have the same objective value at every feasible point,
hence the same endpoints.
\end{proof}

\subsection{Proof of \texorpdfstring{\cref{prop:omitted-bias}}{Proposition on omitted-moderator bias}}
\label{app:proof-omitted-bias}

\begin{proof}
Write $\widetilde C=C-\muC(X)$ and
$\widetilde U=U-\muU(X)$. Randomization and
\eqref{eq:omitted-response-model} imply
\[
  \cov\bigl(C,\pseudo{1}-\pseudo{0}\mid X\bigr)
  =\cov(C\mid X)\theta+\cov(C,U\mid X)\gamma.
\]
Taking expectations gives \eqref{eq:omitted-moment-shift}. Let
$q=\CUcov\gamma$. Every solution of \eqref{eq:omitted-normal-equation} has
the form
\[
  \theta^C=\theta+\CovC^\dagger q+z,
  \qquad z\in\kernel(\CovC),
\]
because $q\in\Vspan=\range(\CovC)$. The true target effect is
\[
  \psi_T=\bT+\ellvec^\top\theta+\ellU^\top\gamma,
\]
which follows by averaging \eqref{eq:omitted-response-model} over the target
covariate law. If $\ct\in\affhull(\pool)$, then
$\ellvec\in\range(\CovC)$ by \cref{lem:loading}, so
$\ellvec^\top z=0$. Subtracting
$\psi_T^C=\bT+\ellvec^\top\theta^C$ gives
\eqref{eq:omitted-bias}.
\end{proof}

\subsection{Proof of \texorpdfstring{\cref{thm:omitted-refinement}}{Theorem on omitted-moderator refinement}}
\label{app:proof-omitted-refinement}

\begin{proof}
Subtract the affine combination of the source conditional treatment
responses from the candidate response. The $g_\Delta(x)$ terms cancel because
$\sum_s w_s=1$, and the context terms cancel because
$\cand{k}=\sum_s w_sc_s$. The remaining term is
\[
  \left(u_k-\sum_s w_su_s\right)^\top\gamma
  =\udisc{k}^\top\gamma,
\]
which proves \eqref{eq:omitted-contrast}.

The candidate's observed contrast therefore adds exactly the equality
$\udisc{k}^\top\gamma=\rho_k$ to the source moment and restriction system,
which proves \eqref{eq:omitted-feasible-update}. The true parameter belongs to
both sets by assumption, so the candidate set is a subset of the source set
and both contain the truth. The infimum of the target functional can only
increase and its supremum can only decrease after intersecting with this
hyperplane, proving \eqref{eq:omitted-monotonicity}. If $\udisc{k}=0$, the
contrast is zero under the model and the added equality is redundant, so the
sets and widths are equal. If the hyperplane cuts a target-relevant direction,
one of the endpoint programs changes strictly.
\end{proof}

\subsection{Proof of \texorpdfstring{\cref{thm:monotone}}{Theorem 1}}
\label{app:proof-monotone}

\begin{proof}
By \cref{prop:refinement}, $\feasset'$ is either $\feasset$ or
$\feasset$ intersected with the two hyperplanes through
$\beta^\circ$, so $\feasset'\subseteq\feasset$ in both cases. The true
vector $\beta^\circ$ satisfies the equality systems of both pools by
\eqref{eq:app-normal} and lies in $\restrset$ by hypothesis, so
$\beta^\circ\in\feasset'$ and both programs are feasible. By
\cref{prop:refinement} the augmented-pool endpoints equal the base
objective optimized over $\feasset'$, so shrinking the feasible set
weakly raises the infimum and weakly lowers the supremum at every
target, which is \eqref{eq:monotone}, with equality throughout in the
rank-preserving case where $\feasset'=\feasset$. Since
$\beta^\circ\in\feasset'$, the value
$\psi_T=\bT+\ellvec^\top(\beta^\circ_1-\beta^\circ_0)$ lies between the
augmented endpoints.
\end{proof}

\subsection{The dual price of the forecast}
\label{app:dual-price}

For completeness we record the sensitivity statement behind
\eqref{eq:dual-pricing}. Fix one polyhedral branch and a target, and
consider the lower endpoint program of \eqref{eq:width-map} as a
function of the added right-hand side $y=(y_0,y_1)$,
\[
  v_-(y)=\bT+\min_{\beta}\bigl\{\ellvec^\top(\beta_1-\beta_0):\
  \CovC\beta_a=\dmom{a},\
  \newdir{k}^\top\beta_a=y_a,\ a\in\{0,1\},\ G\beta\leq h\bigr\}.
\]
This is a linear program whose right-hand side is affine in $y$, so
$v_-$ is a piecewise linear convex function of $y$ on its feasibility
domain, and standard linear-programming duality gives, at any $y$ where
the optimal basis is unique, $\partial v_-/\partial y_a
=\lambda^{u}_{-,a}$, the optimal multiplier on the row
$\newdir{k}^\top\beta_a=y_a$ \citep[Chapter~4]{bonnans2013perturbation}.
The upper endpoint is the corresponding concave value with multipliers
$\lambda^{u}_{+,a}$. At a kink, the left and right slopes are the
extreme multipliers, and the exact range of predicted intervals over
$|y_a-\what y_a|\leq B$ is obtained by re-solving the two programs at
the corners, which is how the simulations and the case study evaluate
forecast uncertainty. For a finite union of branches the value is the
minimum or maximum of finitely many such functions and the same
statements hold branchwise.
